% =============================================================================
% Proxy SOCKS5 - Trabajo Práctico Especial
% Protocolos de Comunicación - ITBA - 2026/1
% =============================================================================
\documentclass[12pt,a4paper]{article}

% ---- Paquetes ----
\usepackage[utf8]{inputenc}
\usepackage[T1]{fontenc}
\usepackage[spanish]{babel}
\usepackage{geometry}
\geometry{left=2.5cm,right=2.5cm,top=2.5cm,bottom=2.5cm}
\usepackage{hyperref}
\usepackage{listings}
\usepackage{xcolor}
\usepackage{graphicx}
\usepackage{enumitem}
\usepackage{float}
\usepackage{fancyhdr}
\usepackage{titlesec}
\usepackage{booktabs}
\usepackage{verbatim}
\usepackage{parskip}

% ---- Configuración de listings ----
\lstset{
    basicstyle=\ttfamily\small,
    breaklines=true,
    frame=single,
    numbers=left,
    numberstyle=\tiny,
    backgroundcolor=\color{gray!10},
    rulecolor=\color{gray!30},
    captionpos=b,
    tabsize=4,
    showstringspaces=false,
    language=C,
    keywordstyle=\color{blue},
    stringstyle=\color{red},
    commentstyle=\color{green!50!black},
    literate=
        {á}{{\'a}}1 {é}{{\'e}}1 {í}{{\'i}}1 {ó}{{\'o}}1 {ú}{{\'u}}1
        {Á}{{\'A}}1 {É}{{\'E}}1 {Í}{{\'I}}1 {Ó}{{\'O}}1 {Ú}{{\'U}}1
        {ñ}{{\~n}}1 {Ñ}{{\~N}}1 {ü}{{\"u}}1 {Ü}{{\"U}}1
        {¿}{{?`}}1 {¡}{{!`}}1
}

% ---- Configuración de encabezados ----
\pagestyle{fancy}
\fancyhf{}
\fancyhead[L]{Proxy SOCKS5}
\fancyhead[R]{Protocolos de Comunicación --- 2026/1}
\fancyfoot[C]{\thepage}

% ---- Redefinir títulos de secciones ----
\titleformat{\section}{\Large\bfseries}{\thesection}{1em}{}
\titleformat{\subsection}{\large\bfseries}{\thesubsection}{1em}{}

% ---- Portada ----
\title{\textbf{Trabajo Práctico Especial}\\[0.5cm]
       \Large{Proxy SOCKS5}\\[0.3cm]
       \large{Protocolos de Comunicación --- 2026/1}}
\author{\textbf{Grupo:} [NOMBRE\_DEL\_GRUPO] \\[0.3cm]
        [Heras\_Tiago] \\[0.1cm]
        [Arias\_Mateo] \\[0.1cm]
        [Vallejo Tapia\_Amador] \\[0.1cm]
        [García Puente\_Manuel] \\[0.3cm]
        \textbf{Instituto Tecnológico de Buenos Aires}}
\date{\today}

\begin{document}

\maketitle
\thispagestyle{empty}
\newpage

% =============================================================================
% 1. ÍNDICE
% =============================================================================
\tableofcontents
\newpage

% =============================================================================
% 2. DESCRIPCIÓN DETALLADA DE LOS PROTOCOLOS DISEÑADOS Y APLICACIONES
% =============================================================================
\section{Descripción detallada de los protocolos diseñados y aplicaciones desarrolladas}

\subsection{Arquitectura general del servidor proxy SOCKS5}

El servidor proxy SOCKS5 está implementado como una aplicación monohilo que utiliza
multiplexación de entrada/salida no bloqueante mediante un selector de eventos
(implementado en \texttt{selector.c}). Todas las conexiones entrantes y salientes
se manejan en un único hilo de ejecución, con la única excepción de la resolución
de nombres de dominio (FQDN), que se delega a un hilo auxiliar mediante
\texttt{getaddrinfo(3)}.

Cada conexión SOCKS5 se modela como una máquina de estados (implementada en
\texttt{stm.c}) que transiciona a través de las siguientes fases:

\begin{verbatim}
HELLO_READ -> HELLO_WRITE -> [AUTH_READ -> AUTH_WRITE] -> REQUEST_READ
    -> REQUEST_RESOLV -> REQUEST_CONNECTING -> REQUEST_WRITE -> COPY
    -> DONE / ERROR
\end{verbatim}

\textbf{Descripción de cada fase:}

\begin{itemize}
    \item \textbf{HELLO}: El cliente envía su versión SOCKS y los métodos de
          autenticación que soporta. El servidor selecciona el método según su
          política interna (sin autenticación si no hay usuarios registrados,
          usuario/contraseña en caso contrario).
    \item \textbf{AUTH (RFC 1929)}: Si se seleccionó el método de usuario/contraseña,
          se intercambian las credenciales y se validan contra la tabla de usuarios
          del servidor.
    \item \textbf{REQUEST}: El cliente envía el pedido de conexión (comando CONNECT,
          tipo de dirección y destino). El servidor parsea la solicitud, resuelve
          el FQDN si es necesario (en un hilo separado), y establece una conexión
          TCP no bloqueante hacia el servidor de origen.
    \item \textbf{COPY}: Una vez establecida la conexión, se inicia un túnel
          bidireccional que copia datos entre el cliente y el servidor de origen
          hasta que alguna de las partes cierra la conexión.
\end{itemize}

\subsection{Protocolo de monitoreo y configuración}

\subsubsection{Descripción general}

El protocolo de monitoreo es un protocolo de texto simple sobre TCP, diseñado
para permitir la administración remota del servidor proxy en tiempo de ejecución
sin necesidad de reiniciar el servicio. Escucha en un socket pasivo independiente
del proxy SOCKS5, en un puerto configurable (por defecto 8080).

\subsubsection{Decisiones de diseño}

Se optó por un protocolo basado en texto plano (en lugar de binario) por las
siguientes razones:

\begin{itemize}
    \item \textbf{Simplicidad de implementación}: Los comandos son legibles por
          humanos, lo que facilita la depuración y las pruebas manuales.
    \item \textbf{Interoperabilidad}: Puede probarse con herramientas estándar
          como \texttt{nc} (netcat) o \texttt{telnet}.
    \item \textbf{Baja complejidad}: Dado que el protocolo maneja pocos comandos
          y respuestas cortas, la sobrecarga del formato texto es despreciable.
    \item \textbf{Extensibilidad}: Agregar nuevos comandos es trivial.
\end{itemize}

\subsubsection{Formato de los mensajes}

\paragraph{Formato de los comandos (cliente $\rightarrow$ servidor)}

Cada comando es una línea de texto ASCII terminada en el carácter \texttt{0x0A}
(\textbackslash n). Opcionalmente puede incluir el carácter \texttt{0x0D}
(\textbackslash r) antes del salto de línea, el cual es ignorado por el servidor.

\begin{lstlisting}[caption=Gramática del protocolo de monitoreo (formato ABNF)]
comando     = verbo *(SP argumento) CRLF
verbo       = 1*ALPHA
argumento   = 1*(%x21-7E)   ; caracteres imprimibles no espacio
CRLF        = [CR] LF
CR          = %x0D
LF          = %x0A
SP          = %x20
\end{lstlisting}

\paragraph{Formato de las respuestas (servidor $\rightarrow$ cliente)}

Cada respuesta es una línea de texto ASCII terminada en \texttt{CRLF}. El primer
token indica éxito o fracaso:

\begin{itemize}
    \item \texttt{+OK} seguido de datos opcionales: el comando se ejecutó
          correctamente.
    \item \texttt{-ERR} seguido de un mensaje de error: el comando falló.
\end{itemize}

\begin{lstlisting}[caption=Gramática de las respuestas]
respuesta    = status [SP datos] CRLF
status       = "+OK" / "-ERR"
datos        = 1*(%x20-7E)
\end{lstlisting}

\subsubsection{Comandos soportados}

\begin{enumerate}
    \item \textbf{GET\_METRICS}
    \begin{itemize}
        \item \textbf{Propósito}: Obtiene las métricas actuales del servidor.
        \item \textbf{Sintaxis}: \texttt{GET\_METRICS\textbackslash n}
        \item \textbf{Respuesta exitosa}: \texttt{+OK ACT:<n> HIST:<n> BYTES:<n>\textbackslash r\textbackslash n}
        \item \textbf{Campos}:
              \begin{itemize}
                  \item \texttt{ACT}: Cantidad de conexiones activas en este momento.
                  \item \texttt{HIST}: Cantidad histórica de conexiones desde el inicio del servidor.
                  \item \texttt{BYTES}: Cantidad total de bytes transferidos.
              \end{itemize}
    \end{itemize}

    \item \textbf{ADD\_USER}
    \begin{itemize}
        \item \textbf{Propósito}: Agrega un nuevo usuario autorizado para usar el proxy.
        \item \textbf{Sintaxis}: \texttt{ADD\_USER <usuario> <contraseña>\textbackslash n}
        \item \textbf{Respuesta exitosa}: \texttt{+OK\textbackslash r\textbackslash n}
        \item \textbf{Respuesta de error}: \texttt{-ERR <motivo>\textbackslash r\textbackslash n}
        \item \textbf{Restricciones}: Máximo 10 usuarios. El nombre de usuario y la
              contraseña no pueden superar los 255 caracteres cada uno.
    \end{itemize}

    \item \textbf{DEL\_USER}
    \begin{itemize}
        \item \textbf{Propósito}: Elimina un usuario previamente registrado.
        \item \textbf{Sintaxis}: \texttt{DEL\_USER <usuario>\textbackslash n}
        \item \textbf{Respuesta exitosa}: \texttt{+OK\textbackslash r\textbackslash n}
        \item \textbf{Respuesta de error}: \texttt{-ERR user not found\textbackslash r\textbackslash n}
    \end{itemize}
\end{enumerate}

\subsubsection{Flujo de una sesión típica}

\begin{lstlisting}[caption=Ejemplo de sesión de monitoreo]
Cliente:                        Servidor:
|  CONECTAR (TCP)               |
|------------------------------>|
|  GET_METRICS\n                |
|------------------------------>|
|                               |  +OK ACT:3 HIST:150 BYTES:1048576\r\n
|<------------------------------|
|  ADD_USER alberto pass123\n   |
|------------------------------>|
|                               |  +OK\r\n
|<------------------------------|
|  DEL_USER alberto\n           |
|------------------------------>|
|                               |  +OK\r\n
|<------------------------------|
|  CERRAR CONEXIÓN              |
|------------------------------>|
\end{lstlisting}

\subsection{Aplicación cliente de monitoreo}

El cliente de monitoreo es una aplicación de línea de comandos que utiliza I/O
bloqueante (dada la simplicidad del protocolo). Se conecta al servidor de
management, envía el comando solicitado y muestra la respuesta.

\subsubsection{Sintaxis}

\begin{lstlisting}[caption=Uso del cliente de monitoreo]
./bin/client [opciones] <comando>

Opciones:
  -L <dirección>   Dirección IP del servidor de management (defecto: 127.0.0.1)
  -P <puerto>      Puerto del servidor de management (defecto: 8080)

Comandos:
  get-metrics                  Obtiene métricas del servidor
  add-user <usr> <pass>        Agrega un usuario
  del-user <usr>               Elimina un usuario
\end{lstlisting}

\subsubsection{Ejemplos de uso}

\begin{lstlisting}[caption=Ejemplos del cliente de monitoreo]
$ ./bin/client get-metrics
+OK ACT:0 HIST:12 BYTES:409600

$ ./bin/client add-user juan pass123
+OK

$ ./bin/client del-user juan
+OK
\end{lstlisting}

\subsection{Métricas del sistema}

El servidor recolecta las siguientes métricas, accesibles a través del comando
\texttt{GET\_METRICS}:

\begin{itemize}
    \item \textbf{Conexiones activas} (\texttt{ACT}): Número de conexiones SOCKS5
          actualmente en curso. Se incrementa al aceptar una nueva conexión y se
          decrementa al finalizar.
    \item \textbf{Conexiones históricas} (\texttt{HIST}): Contador acumulativo de
          todas las conexiones SOCKS5 establecidas desde el inicio del servidor.
    \item \textbf{Bytes transferidos} (\texttt{BYTES}): Suma total de bytes
          transferidos en ambas direcciones (cliente $\rightarrow$ proxy y
          proxy $\rightarrow$ origen) a través del túnel.
\end{itemize}

Todas las métricas son volátiles: se pierden al reiniciar el servidor.

% =============================================================================
% 3. PROBLEMAS ENCONTRADOS
% =============================================================================
\section{Problemas encontrados durante el diseño y la implementación}

\subsection{[Aca tenemos que poner los problemas que nos surgieron]}

\begin{itemize}
    \item \textbf{Manejo de lecturas/escrituras parciales}: Durante la fase de
          COPY, fue necesario implementar un mecanismo robusto que maneje
          correctamente las lecturas y escrituras parciales, propagando los
          cierres de conexión de forma adecuada entre ambas mitades del túnel.
    \item \textbf{Resolución de nombres no bloqueante}: Dado que
          \texttt{getaddrinfo(3)} es una operación bloqueante, se optó por
          ejecutarla en un hilo separado y notificar al selector mediante
          \texttt{selector\_notify\_block}. Esto requirió cuidado para evitar
          condiciones de carrera en el acceso a las estructuras compartidas.
    \item \textbf{[Terminar con cosas que tuvimos problemas]}
\end{itemize}

% =============================================================================
% 4. LIMITACIONES
% =============================================================================
\section{Limitaciones de la aplicación}

\subsection{Registro de accesos bloqueante}

El registro de accesos (\texttt{log\_access}) escribe en el archivo
\texttt{access.log} utilizando operaciones de E/S bloqueantes
(\texttt{fopen}, \texttt{fprintf}, \texttt{fclose}) dentro del bucle principal
de eventos no bloqueante. Esto puede introducir micro-pausas que degradan el
rendimiento bajo carga elevada. Una mejora posible sería implementar un buffer
de logs en memoria con escritura asincrónica.

\subsection{Autenticación del protocolo de monitoreo}

El servidor de management no implementa autenticación. Cualquier entidad que
pueda conectarse al puerto de monitoreo (por defecto 8080 en \texttt{127.0.0.1})
puede agregar o eliminar usuarios del proxy. Si bien el puerto de monitoreo está
configurado por defecto para escuchar solo en la interfaz loopback, esta sigue
siendo una limitación de seguridad.

\subsection{Backlog de conexiones}

El socket pasivo del proxy SOCKS5 se configura con un backlog de 20 conexiones
(\texttt{listen(server, 20)}). Para soportar 500 clientes concurrentes de forma
robusta, sería recomendable utilizar \texttt{SOMAXCONN} (generalmente 128 o
superior).

\subsection{Soporte limitado de comandos SOCKS}

Actualmente solo se implementa el comando \texttt{CONNECT}. Los comandos
\texttt{BIND} y \texttt{UDP ASSOCIATE} definidos en RFC 1928 no están
soportados.

\subsection{[Aca podemos agregar mas limitaciones y quizas podemos arreglar algunas de las de arriba]}

% =============================================================================
% 5. POSIBLES EXTENSIONES
% =============================================================================
\section{Posibles extensiones}

\begin{itemize}
    \item \textbf{Dissector de credenciales (ettercap-like)}: Implementar un
          módulo de inspección de tráfico en la fase COPY que detecte
          credenciales enviadas en texto plano en protocolos como POP3, FTP,
          HTTP Basic, SMTP, etc. Esta funcionalidad está prevista para la
          segunda entrega del trabajo práctico.
    \item \textbf{Soporte BIND y UDP ASSOCIATE}: Completar la implementación
          de los comandos faltantes del RFC 1928.
    \item \textbf{Autenticación del protocolo de monitoreo}: Agregar un
          mecanismo de autenticación (contraseña compartida, token, o
          autenticación mutua) para el servidor de management.
    \item \textbf{Logging asincrónico}: Reemplazar el registro de accesos
          bloqueante por un esquema con buffer circular y escritura diferida.
    \item \textbf{Configuración dinámica extendida}: Permitir cambiar en
          tiempo de ejecución otros parámetros como timeouts, tamaño de
          buffers, límite de conexiones concurrentes, etc.
    \item \textbf{Estadísticas avanzadas}: Agregar métricas como tiempo de
          vida promedio de conexiones, throughput por usuario, histogramas de
          tamaño de transferencias, etc.
\end{itemize}

% =============================================================================
% 6. CONCLUSIONES
% =============================================================================
\section{Conclusiones}

[Aca pongamos unas conclusiones entre todos si quieren]

% =============================================================================
% 7. EJEMPLOS DE PRUEBA
% =============================================================================
\section{Ejemplos de prueba}

\subsection{Prueba básica del proxy SOCKS5}

\begin{lstlisting}[caption=Iniciar el servidor proxy]
# Iniciar el servidor con un usuario
$ ./bin/server -u testuser:testpass
Listening on TCP port 1080
Management listening on TCP port 8080
\end{lstlisting}

\begin{lstlisting}[caption=Probar el proxy con curl]
# Solicitud HTTP a través del proxy SOCKS5 con autenticación
$ curl --socks5 127.0.0.1:1080 --proxy-user testuser:testpass \
       http://example.com
<!doctype html>
<html>
...
\end{lstlisting}

\begin{lstlisting}[caption=Probar el proxy sin autenticación]
# Iniciar sin usuarios registrados (acepta "sin autenticación")
$ ./bin/server
$ curl --socks5 127.0.0.1:1080 http://example.com
\end{lstlisting}

\subsection{Prueba de monitoreo}

\begin{lstlisting}[caption=Consultar métricas]
$ ./bin/client get-metrics
+OK ACT:1 HIST:5 BYTES:24576
\end{lstlisting}

\begin{lstlisting}[caption=Agregar y eliminar usuarios en caliente]
$ ./bin/client add-user nuevo pass123
+OK
$ ./bin/client del-user nuevo
+OK
\end{lstlisting}

\subsection{Prueba de graceful shutdown}

\begin{lstlisting}[caption=Enviar señal SIGTERM al servidor]
$ kill -TERM <PID_del_servidor>
# El servidor imprime:
# graceful shutdown: no se aceptan nuevas conexiones, esperando N conexiones activas
# (espera a que terminen todas las conexiones activas y luego finaliza)
\end{lstlisting}

% =============================================================================
% 8. GUÍA DE INSTALACIÓN
% =============================================================================
\section{Guía de instalación}

\subsection{Requisitos del sistema}

\begin{itemize}
    \item Sistema operativo: Linux (kernel 4.x o superior)
    \item Compilador: GCC 4.9 o superior (compatible con C11)
    \item GNU Make 3.81 o superior
    \item Bibliotecas: libc estándar, pthreads
\end{itemize}

\subsection{Compilación}

\begin{lstlisting}[caption=Compilar el proyecto desde la raíz]
$ make clean && make all
\end{lstlisting}

Este comando genera los siguientes artefactos:

\begin{itemize}
    \item \texttt{bin/server} --- Servidor proxy SOCKS5
    \item \texttt{bin/client} --- Cliente de monitoreo
\end{itemize}

Los archivos objeto intermedios se almacenan en el directorio \texttt{obj/}.

\subsection{Limpieza}

\begin{lstlisting}[caption=Eliminar artefactos compilados]
$ make clean
\end{lstlisting}

% =============================================================================
% 9. INSTRUCCIONES DE CONFIGURACIÓN
% =============================================================================
\section{Instrucciones para la configuración}

\subsection{Opciones del servidor}

\begin{lstlisting}[caption=Uso del servidor]
./bin/server [OPCIONES]

Opciones:
  -h               Imprime la ayuda y termina.
  -l <dirección>   Dirección donde servirá el proxy SOCKS (defecto: 0.0.0.0).
  -L <dirección>   Dirección donde servirá el servicio de management
                   (defecto: 127.0.0.1).
  -p <puerto>      Puerto entrante para conexiones SOCKS (defecto: 1080).
  -P <puerto>      Puerto entrante para conexiones de management
                   (defecto: 8080).
  -u <usr>:<pass>  Usuario y contraseña para el proxy. Hasta 10 ocurrencias.
  -v               Imprime la versión y termina.
\end{lstlisting}

\subsection{Ejemplos de configuración}

\begin{lstlisting}[caption=Servidor con configuración personalizada]
# Servidor SOCKS en puerto 2080, management en puerto 9090
$ ./bin/server -p 2080 -P 9090

# Servidor con dos usuarios predefinidos
$ ./bin/server -u alice:pass1 -u bob:pass2

# Servidor con management en una interfaz específica
$ ./bin/server -L 192.168.1.100 -P 9090
\end{lstlisting}

\subsection{Opciones del cliente de monitoreo}

\begin{lstlisting}[caption=Uso del cliente]
./bin/client [OPCIONES] <comando>

Opciones:
  -L <dirección>   Dirección IP del servidor de management (defecto: 127.0.0.1)
  -P <puerto>      Puerto del servidor de management (defecto: 8080)

Comandos:
  get-metrics                  Obtiene métricas del servidor
  add-user <usr> <pass>        Agrega un usuario
  del-user <usr>               Elimina un usuario
\end{lstlisting}

% =============================================================================
% 10. EJEMPLOS DE CONFIGURACIÓN Y MONITOREO
% =============================================================================
\section{Ejemplos de configuración y monitoreo}

\subsection{Escenario 1: Servidor básico con monitoreo local}

\begin{lstlisting}[caption=Inicialización y monitoreo]
# Terminal 1: Iniciar el servidor
$ ./bin/server -u admin:admin123

# Terminal 2: Consultar métricas iniciales
$ ./bin/client get-metrics
+OK ACT:0 HIST:0 BYTES:0

# Terminal 2: Agregar un usuario nuevo
$ ./bin/client add-user desarrollador pass123
+OK

# Terminal 2: Verificar que el usuario se agregó (conectándose al proxy)
$ curl --socks5 127.0.0.1:1080 --proxy-user desarrollador:pass123 \
       https://httpbin.org/ip

# Terminal 2: Consultar métricas después de la conexión
$ ./bin/client get-metrics
+OK ACT:0 HIST:1 BYTES:2048
\end{lstlisting}

\subsection{Escenario 2: Acceso a un sitio bloqueado}

\begin{lstlisting}[caption=Verificación del registro de accesos]
# Después de varias conexiones, revisar el archivo de log
$ cat access.log
[2026-07-06 20:15:30] User: admin | FQDN: example.com | IP Destino: 93.184.216.34:80
[2026-07-06 20:16:45] User: admin | FQDN: httpbin.org | IP Destino: 3.220.72.84:443
\end{lstlisting}

% =============================================================================
% 11. DOCUMENTO DE DISEÑO
% =============================================================================
\section{Documento de diseño del proyecto}

\subsection{Diagrama de módulos}

\begin{verbatim}
+------------------------------------------------------------------+
|                        bin/server                                 |
|                                                                   |
|  +------------------+  +------------------+  +-----------------+  |
|  |   main.c         |  |   args.c/h       |  |  mng_server.c/h |  |
|  |  - bucle ppal    |  |  - parseo CLI    |  |  - protocolo    |  |
|  |  - graceful      |  |  - defaults      |  |    monitoreo    |  |
|  |    shutdown      |  +------------------+  |  - ADD/DEL_USER |  |
|  +--------+---------+                       |  - GET_METRICS   |  |
|           |                                  +--------+--------+  |
|           v                                           |           |
|  +------------------+                                  |           |
|  |  socks5nio.c/h   |<--+                             |           |
|  |  - stm SOCKS5    |   |                             |           |
|  |  - HELLO/AUTH/   |   |                             |           |
|  |    REQUEST/COPY  |   |                             |           |
|  +--------+---------+   |                             |           |
|           |             |                             |           |
|  +--------+---------+   |                             |           |
|  |  hello.c/h       |   |  parsers                   |           |
|  |  auth.c/h        |---+  SOCKS5                    |           |
|  |  request.c/h     |   |                            |           |
|  +------------------+   |                            |           |
|                         |                            |           |
|  +------------------+   |  +------------------+      |           |
|  |  selector.c/h    |<--+--|  stm.c/h         |      |           |
|  |  - multiplexor   |      |  - máquina de    |      |           |
|  |    I/O no blq.   |      |    estados       |      |           |
|  +------------------+      +------------------+      |           |
|                                                                   |
|  +------------------+  +------------------+                       |
|  |  buffer.c/h      |  |  netutils.c/h    |                       |
|  |  - buffer circ.  |  |  - utilidades    |                       |
|  |    con acceso     |  |    de red        |                       |
|  |    directo        |  |                  |                       |
|  +------------------+  +------------------+                       |
+------------------------------------------------------------------+

+------------------------------------------------------------------+
|                        bin/client                                 |
|                                                                   |
|  +------------------+                                             |
|  |  client.c        |                                             |
|  |  - I/O bloqueante|                                             |
|  |  - envía comando |                                             |
|  |  - muestra resp. |                                             |
|  +------------------+                                             |
+------------------------------------------------------------------+
\end{verbatim}

\subsection{Descripción de los componentes}

\subsubsection{Selector de eventos (\texttt{selector.c/h})}

El selector es el núcleo del multiplexado de E/S. Permite registrar
descriptores de archivo con intereses de lectura/escritura y notifica
mediante callbacks cuando los eventos ocurren. Soporta notificación de
operaciones bloqueantes (como resolución de DNS) mediante señales.

\textbf{Interfaz pública:}
\begin{itemize}
    \item \texttt{selector\_init()}: Inicializa la librería con una señal
          para notificaciones internas.
    \item \texttt{selector\_new(cantidad)}: Crea una nueva instancia de
          selector con capacidad para \texttt{cantidad} descriptores.
    \item \texttt{selector\_register()}: Registra un fd con un handler y
          un interés inicial.
    \item \texttt{selector\_select()}: Espera eventos y dispara los
          callbacks correspondientes.
    \item \texttt{selector\_notify\_block()}: Notifica al selector que
          una operación bloqueante (ej: DNS) finalizó.
\end{itemize}

\subsubsection{Buffer circular (\texttt{buffer.c/h})}

Buffer con acceso directo a memoria que mantiene punteros de lectura y
escritura separados. Permite operaciones de I/O eficientes al exponer
directamente las direcciones de memoria donde leer o escribir, evitando
copias intermedias.

\textbf{Características:}
\begin{itemize}
    \item Compactación automática cuando los punteros se igualan.
    \item Sin asignaciones dinámicas de memoria (usa buffers pre-asignados).
    \item Ideal para operaciones de \texttt{recv(2)} y \texttt{send(2)}
          no bloqueantes.
\end{itemize}

\subsubsection{Máquina de estados (\texttt{stm.c/h})}

Motor de máquina de estados genérico donde los eventos son los del
selector (\texttt{on\_read\_ready}, \texttt{on\_write\_ready},
\texttt{on\_block\_ready}). Cada estado define callbacks de llegada,
salida y manejo de eventos.

\textbf{Estados definidos para SOCKS5:}
\begin{itemize}
    \item \texttt{HELLO\_READ}, \texttt{HELLO\_WRITE}
    \item \texttt{AUTH\_READ}, \texttt{AUTH\_WRITE}
    \item \texttt{REQUEST\_READ}, \texttt{REQUEST\_RESOLV},
          \texttt{REQUEST\_CONNECTING}, \texttt{REQUEST\_WRITE}
    \item \texttt{COPY}
    \item \texttt{DONE}, \texttt{ERROR}
\end{itemize}

\subsubsection{Parsers SOCKS5}

\begin{itemize}
    \item \textbf{hello.c/h}: Parser del mensaje HELLO (RFC 1928, sección 3).
          Procesa versión, cantidad de métodos y lista de métodos de
          autenticación. Incluye callback por método para selección dinámica.
    \item \textbf{auth.c/h}: Parser de autenticación usuario/contraseña
          (RFC 1929). Procesa los campos VER, ULEN, UNAME, PLEN, PASSWD.
    \item \textbf{request.c/h}: Parser del mensaje REQUEST (RFC 1928,
          sección 4). Soporta direcciones IPv4, IPv6 y FQDN. Incluye
          serialización de respuestas con BND.ADDR.
\end{itemize}

\subsubsection{Servidor de monitoreo (\texttt{mng\_server.c/h})}

Implementa el protocolo de monitoreo como una máquina de estados simple
(\texttt{MNG\_READ} $\rightarrow$ \texttt{MNG\_EXECUTE} $\rightarrow$
\texttt{MNG\_WRITE} $\rightarrow$ \texttt{MNG\_DONE}). Utiliza el mismo
selector de eventos que el proxy SOCKS5, integrándose sin necesidad de
hilos adicionales.

\subsubsection{Cliente de monitoreo (\texttt{client.c})}

Aplicación de línea de comandos con I/O bloqueante. Conecta por TCP al
servidor de management, envía el comando y muestra la respuesta. Soporta
los comandos \texttt{get-metrics}, \texttt{add-user} y \texttt{del-user}.

\subsection{Diagrama de flujo de una conexión SOCKS5}

\begin{verbatim}
Cliente SOCKS5              Servidor Proxy              Servidor Origen
      |                          |                            |
      |--- HELLO (VER, NMETHODS) |                            |
      |------------------------->|                            |
      |                          |--- (selecciona método)     |
      |<---- HELLO (VER, METHOD) |                            |
      |                          |                            |
      | (si METHOD = 0x02)       |                            |
      |--- AUTH (VER, ULEN, ...) |                            |
      |------------------------->|                            |
      |                          |--- (valida credenciales)   |
      |<---- AUTH (VER, STATUS)  |                            |
      |                          |                            |
      |--- REQUEST (CMD, ATYP,..)|                            |
      |------------------------->|                            |
      |                          |--- (resuelve FQDN si      |
      |                          |     es necesario)          |
      |                          |                            |
      |                          |======= CONNECT ===========>|
      |                          |<======= SYN-ACK ===========|
      |                          |                            |
      |<--- REQUEST (REP, BND)   |                            |
      |                          |                            |
      |=================== TÚNEL BIDIRECCIONAL ==============>|
      |<====================================================>|
      |                          |                            |
\end{verbatim}

% =============================================================================
% FIN DEL DOCUMENTO
% =============================================================================
\end{document}